\chapter{Physical and Mathematical Formulation}
\label{chap:physical_mathematical_formulation}

%======================================================================
\section{Physical Description of Hydrogen Transport}
\label{sec:physical_description}
%======================================================================

\subsection{Physical Domain and Dimensionality}
\label{subsec:physical_domain}

The physical problem considered in this thesis involves non-isothermal,
transient hydrogen transport, microstructural trapping and detrapping,
and thermal desorption in ferritic pipeline steel. The specimen is
represented as an isotropic, homogeneous one-dimensional continuum slab
of total thickness $L$, or equivalently as a half-domain of thickness
$a=L/2$ when geometric symmetry is exploited
\cite{Turnbull1997,Legrand2015,Diaz2020,Marrani2025}.

The spatial domain is therefore defined along the thickness coordinate as

\begin{equation}
\Omega = \left[0,\frac{L}{2}\right],
\qquad
x \in \left[0,\frac{L}{2}\right].
\label{eq:ch3_domain}
\end{equation}

Here, $x=0$ represents the mid-plane of the specimen and
$x=L/2$ represents the external surface through which hydrogen exits
during thermal desorption. The mid-plane is treated as a plane of
symmetry, while the external surface is the hydrogen desorption
boundary.

The one-dimensional approximation is appropriate when the specimen
thickness is considerably smaller than its lateral dimensions, such
that the dominant concentration gradients and associated hydrogen
transport occur through the thickness direction
\cite{Turnbull1997,Legrand2015,Marrani2025}. Consequently, lateral
variations in hydrogen concentration are neglected in the present
formulation.

The model assumes that the specimen temperature is spatially uniform
at each instant of time. Thus, although the problem is transient and
non-isothermal, temperature is prescribed as a function of time rather
than solved from a separate heat-conduction equation. The temperature
field is therefore represented by

\begin{equation}
T=T(t).
\label{eq:ch3_uniform_temperature}
\end{equation}

%----------------------------------------------------------------------
\subsection{Hydrogen Species and Concentration Variables}
\label{subsec:hydrogen_species}
%----------------------------------------------------------------------

Hydrogen within the ferritic steel matrix is separated into mobile
lattice hydrogen and hydrogen trapped at microstructural defects
\cite{Oriani1970,Turnbull1997,Krom2000,Legrand2015}.

\paragraph{Interstitial lattice hydrogen.}

Interstitial lattice hydrogen denotes hydrogen atoms occupying normal
interstitial lattice sites (NILS) of the body-centered cubic (BCC)
$\alpha$-iron lattice. This population is considered mobile and is
responsible for spatial diffusion through the steel matrix
\cite{Oriani1970,Chen2025,Fallahmohammadi2013,Legrand2015}.

The lattice hydrogen concentration is denoted by

\begin{equation}
C_L=C_L(x,t),
\label{eq:ch3_CL_definition}
\end{equation}

where $C_L represents the molar concentration of mobile lattice
hydrogen, with units of

\begin{equation}
[C_L]=\mathrm{mol\,m^{-3}}.
\end{equation}

The corresponding density of normal interstitial lattice sites is
denoted by $N_L$, also expressed on a molar-volume basis in the present
formulation.

\paragraph{Trapped hydrogen.}

Hydrogen may also be captured by microstructural defects that act as
energetically favourable trapping sites. Examples include
dislocations, vacancies, grain boundaries, interfaces, and other
microstructural features capable of producing local potential-energy
wells \cite{Darken1949,McNabb1963,Oriani1970,Krom2000}.

The material is represented by $n_t$ discrete and non-interacting trap
populations. The concentration of hydrogen associated with the
$i$-th trap population is written as

\begin{equation}
C_T^{(i)}=C_T^{(i)}(x,t),
\qquad
i=1,2,\ldots,n_t,
\label{eq:ch3_CT_definition}
\end{equation}

with

\begin{equation}
[C_T^{(i)}]=\mathrm{mol\,m^{-3}}.
\end{equation}

The total hydrogen concentration at a given location is consequently

\begin{equation}
C_{\mathrm{H}}(x,t)
=
C_L(x,t)
+
\sum_{i=1}^{n_t} C_T^{(i)}(x,t).
\label{eq:ch3_total_hydrogen}
\end{equation}

The distinction between lattice and trapped populations is important
because only the lattice population is directly transported by the
Fickian diffusion flux in the present formulation, whereas trapped
hydrogen exchanges with the lattice through local trapping and
detrapping reactions.

\paragraph{Lattice occupancy.}

The fractional occupancy of normal interstitial lattice sites is
defined as

\begin{equation}
\theta_L(x,t)
=
\frac{C_L(x,t)}{N_L}.
\label{eq:ch3_theta_L}
\end{equation}

Similarly, the fractional occupancy of the $i$-th trap population is

\begin{equation}
\theta_T^{(i)}(x,t)
=
\frac{C_T^{(i)}(x,t)}{N_T^{(i)}},
\qquad
i=1,2,\ldots,n_t,
\label{eq:ch3_theta_T}
\end{equation}

where $N_T^{(i)}$ is the total density of available sites belonging to
the $i$-th trap population.

These occupancy variables provide a convenient representation of the
finite availability of both lattice and trap sites and are used in the
McNabb--Foster trapping kinetics
\cite{Krom2000,Legrand2015,Marrani2025}.

%----------------------------------------------------------------------
\subsection{Model Identity and Trapping Framework}
\label{subsec:model_identity}
%----------------------------------------------------------------------

The present formulation is a one-dimensional, multi-trap,
non-isothermal transient hydrogen transport model based on
McNabb--Foster kinetic trapping and detrapping, with an Oriani
local-equilibrium condition used to establish the initial hydrogen
distribution \cite{Oriani1970,McNabb1963,Turnbull1997,Krom2000,
Diaz2020,Marrani2025}.

The two formulations are assigned distinct roles.

\begin{itemize}

    \item \textbf{Oriani local equilibrium} is used to determine the
    initial partitioning of hydrogen between lattice and trap
    populations at the initial temperature $T_0$.

    \item \textbf{McNabb--Foster kinetics} governs the subsequent
    time-dependent exchange of hydrogen between the lattice and the
    trap populations during the thermal history, including both
    trapping and detrapping processes.

\end{itemize}

Thus, the present model does not impose instantaneous local
equilibrium between lattice and trapped hydrogen throughout the
thermal desorption experiment. Instead, finite trapping and detrapping
rates are explicitly retained. Consequently, during a finite-rate
temperature ramp, the system can depart from instantaneous local
equilibrium and exhibit transient kinetic effects
\cite{Turnbull1997,Legrand2015,Diaz2020,Marrani2025}.

In the limiting case in which the kinetic exchange becomes
sufficiently rapid, the McNabb--Foster formulation approaches the
corresponding local-equilibrium behaviour represented by the Oriani
isotherm \cite{Turnbull1997,Krom2000}.

%----------------------------------------------------------------------
\subsection{Fundamental Physical Assumptions}
\label{subsec:physical_assumptions}
%----------------------------------------------------------------------

The following assumptions define the physical scope of the present
hydrogen transport and thermal desorption model:

\begin{enumerate}

    \item \textbf{One-dimensional transport:}
    hydrogen transport occurs only along the specimen thickness
    coordinate $x$.

    \item \textbf{Isotropic and homogeneous matrix:}
    the ferritic steel matrix is treated as a continuum with spatially
    uniform material properties, unless a temperature dependence is
    explicitly introduced.

    \item \textbf{Fixed trap population:}
    the number density of trap sites $N_T^{(i)}$ remains constant
    throughout the simulation,

    \begin{equation}
    \frac{\partial N_T^{(i)}}{\partial t}=0.
    \label{eq:ch3_constant_trap_density}
    \end{equation}

    Trap generation, annihilation, and microstructural phase
    transformations are therefore not considered.

    \item \textbf{Non-interacting trap populations:}
    the individual trap populations are treated as independent
    populations. Direct hydrogen transfer from one trap population to
    another is neglected; hydrogen exchange occurs through the mobile
    lattice population.

    \item \textbf{Fickian lattice diffusion:}
    only the mobile lattice population contributes directly to the
    spatial diffusion flux.

    \item \textbf{Stress-free desorption formulation:}
    stress-assisted hydrogen drift caused by hydrostatic stress
    gradients is neglected during the thermal desorption calculation.

    \item \textbf{Uniform instantaneous temperature:}
    the specimen is assumed to have a spatially uniform temperature
    $T(t)$ at each instant.

    \item \textbf{No independent thermal field solution:}
    the temperature history is prescribed through the heating
    schedule rather than being obtained from a coupled heat-transfer
    calculation.

\end{enumerate}

%======================================================================
\section{Hydrogen Diffusion}
\label{sec:hydrogen_diffusion}
%======================================================================

\subsection{Fickian Lattice Flux}
\label{subsec:fickian_flux}

The spatial transport of mobile interstitial hydrogen through the
normal lattice sites is described by Fick's first law
\cite{Fallahmohammadi2013,Legrand2015,MartinezPaneda2020}:

\begin{equation}
J_L(x,t)
=
-D_L(T)
\frac{\partial C_L(x,t)}{\partial x},
\label{eq:ch3_fick_flux}
\end{equation}

where $J_L$ is the lattice hydrogen flux and $D_L(T)$ is the
temperature-dependent lattice diffusion coefficient.

The units are

\begin{equation}
[J_L]=\mathrm{mol\,m^{-2}\,s^{-1}},
\end{equation}

and

\begin{equation}
[D_L]=\mathrm{m^2\,s^{-1}}.
\end{equation}

The negative sign indicates that hydrogen diffuses from regions of
higher lattice concentration toward regions of lower lattice
concentration.

%----------------------------------------------------------------------
\subsection{Temperature-Dependent Lattice Diffusivity}
\label{subsec:lattice_diffusivity}
%----------------------------------------------------------------------

The lattice diffusion coefficient is described by an Arrhenius
temperature dependence:

\begin{equation}
D_L(T)
=
D_0
\exp\left(
-\frac{E_L}{RT(t)}
\right),
\label{eq:ch3_arrhenius_DL}
\end{equation}

where $D_0$ is the pre-exponential diffusion coefficient and $E_L$ is
the activation energy for lattice diffusion. Their units are

\begin{equation}
[D_0]=\mathrm{m^2\,s^{-1}},
\qquad
[E_L]=\mathrm{J\,mol^{-1}}.
\end{equation}

The universal gas constant is

\begin{equation}
R=8.31446\,
\mathrm{J\,mol^{-1}\,K^{-1}},
\label{eq:ch3_R}
\end{equation}

and $T(t)$ is the absolute temperature in kelvin.

The temperature dependence of $D_L$ provides the coupling between the
thermal history and the rate of lattice hydrogen transport.

%----------------------------------------------------------------------
\subsection{McNabb--Foster Transient Trapping Kinetics}
\label{subsec:mcnabb_foster}
%----------------------------------------------------------------------

The exchange of hydrogen between the mobile lattice and each trap
population is described using the McNabb--Foster kinetic formulation
\cite{McNabb1963,Turnbull1997,Krom2000,Legrand2015,Marrani2025}.

For the $i$-th trap population,

\begin{equation}
\frac{\partial C_T^{(i)}}{\partial t}
=
j_{\mathrm{trap}}^{(i)}
=
k_i(T)
\frac{N_T^{(i)}}{N_L}
C_L
\left(
1-\frac{C_T^{(i)}}{N_T^{(i)}}
\right)
-
p_i(T)
C_T^{(i)}
\left(
1-\frac{C_L}{N_L}
\right),
\label{eq:ch3_mcnabb_molar}
\end{equation}

where $k_i(T)$ is the trapping rate coefficient and $p_i(T)$ is the
detrapping rate coefficient.

The first term represents hydrogen transfer from the mobile lattice
into vacant trap sites. The second term represents release of hydrogen
from occupied trap sites back into available lattice sites.

Using the fractional occupancies defined previously, the same
equation can be written as

\begin{equation}
\frac{\partial \theta_T^{(i)}}{\partial t}
=
k_i(T)\theta_L
\left(1-\theta_T^{(i)}\right)
-
p_i(T)\theta_T^{(i)}
\left(1-\theta_L\right).
\label{eq:ch3_mcnabb_occupancy}
\end{equation}

The trapping contribution increases with lattice occupancy
$\theta_L$ and available trap capacity
$(1-\theta_T^{(i)})$, whereas the detrapping contribution increases
with trap occupancy $\theta_T^{(i)}$ and available lattice capacity
$(1-\theta_L)$.

%----------------------------------------------------------------------
\subsection{Temperature Dependence of Trapping and Detrapping}
\label{subsec:trapping_detrapping_arrhenius}
%----------------------------------------------------------------------

The trapping and detrapping rate coefficients are represented using
Arrhenius-type expressions based on the lattice hopping activation
energy and the trap binding energy
\cite{Turnbull1997,Legrand2015,Diaz2020,Marrani2025}.

The trapping rate coefficient is

\begin{equation}
k_i(T)
=
\nu_t
\exp\left(
-\frac{E_L}{RT(t)}
\right),
\label{eq:ch3_arrhenius_ki}
\end{equation}

while the detrapping rate coefficient is

\begin{equation}
p_i(T)
=
\nu_d
\exp\left(
-\frac{E_L+E_b^{(i)}}{RT(t)}
\right)
=
\nu_d
\exp\left(
-\frac{E_d^{(i)}}{RT(t)}
\right).
\label{eq:ch3_arrhenius_pi}
\end{equation}

Here, $\nu_t$ and $\nu_d$ are the attempt frequencies associated with
trapping and detrapping, respectively. In the present formulation,
they may be taken as equal to the characteristic lattice jump or
Debye frequency,

\begin{equation}
\nu_t=\nu_d=\nu
\approx 10^{13}\,\mathrm{s^{-1}},
\label{eq:ch3_debye_frequency}
\end{equation}

following the formulation adopted in the cited ferritic-iron
transport models \cite{Turnbull1997,Legrand2015,Diaz2020,Marrani2025}.

The quantity $E_b^{(i)}$ represents the positive binding energy of the
$i$-th trap population. With the convention used here,

\begin{equation}
E_b^{(i)}=-\Delta H_i,
\qquad
E_b^{(i)}>0
\quad\text{for an attractive trap}.
\label{eq:ch3_binding_energy_convention}
\end{equation}

The corresponding detrapping activation energy is

\begin{equation}
E_d^{(i)}
=
E_L+E_b^{(i)}.
\label{eq:ch3_detrapping_energy}
\end{equation}

Thus, deeper traps possess larger detrapping activation energies and
require a larger thermal driving force for hydrogen release.

%----------------------------------------------------------------------
\subsection{Combined Governing Transient Conservation Equation}
\label{subsec:governing_pde}
%----------------------------------------------------------------------

Conservation of hydrogen requires that the temporal change in the
total hydrogen concentration within a differential control volume is
balanced by the divergence of the mobile lattice diffusion flux
\cite{McNabb1963,Turnbull1997,Legrand2015,Marrani2025}.

The resulting governing equation is

\begin{equation}
\frac{\partial C_L(x,t)}{\partial t}
+
\sum_{i=1}^{n_t}
\frac{\partial C_T^{(i)}(x,t)}{\partial t}
=
\frac{\partial}{\partial x}
\left[
D_L(T)
\frac{\partial C_L(x,t)}{\partial x}
\right].
\label{eq:ch3_governing_pde}
\end{equation}

The equation contains two coupled physical processes:

\begin{enumerate}
    \item spatial diffusion of mobile lattice hydrogen, represented by
    the right-hand side; and
    \item local exchange between lattice and trapped populations,
    represented by the temporal derivatives of the trapped
    concentrations.
\end{enumerate}

Substitution of the McNabb--Foster trapping equations into the
conservation equation gives the explicit nonlinear diffusion-reaction
equation

\begin{equation}
\begin{aligned}
\frac{\partial C_L}{\partial t}
&+
\sum_{i=1}^{n_t}
\left[
k_i
\frac{N_T^{(i)}}{N_L}
C_L
\left(
1-\frac{C_T^{(i)}}{N_T^{(i)}}
\right)
-
p_i C_T^{(i)}
\left(
1-\frac{C_L}{N_L}
\right)
\right]
\\
&=
\frac{\partial}{\partial x}
\left[
D_L(T)
\frac{\partial C_L}{\partial x}
\right].
\end{aligned}
\label{eq:ch3_governing_pde_expanded}
\end{equation}

This equation constitutes the primary physical governing equation
solved by the numerical model. The nonlinearities arise from the
finite lattice and trap occupancies and from the temperature
dependence of the diffusion, trapping, and detrapping coefficients.

%----------------------------------------------------------------------
\subsection{Initial Conditions and Oriani Equilibrium Initialization}
\label{subsec:initial_conditions}
%----------------------------------------------------------------------

At the beginning of the thermal desorption calculation, the mobile
lattice hydrogen concentration is prescribed as a spatially uniform
initial value:

\begin{equation}
C_L(x,0)=C_{L,0},
\qquad
0\leq x\leq \frac{L}{2}.
\label{eq:ch3_ic_CL}
\end{equation}

The initial trapped hydrogen concentrations are then determined from
the Oriani local-equilibrium relation at the initial temperature
$T_0$ \cite{Oriani1970,Diaz2020,Marrani2025}.

For the $i$-th trap population,

\begin{equation}
\theta_T^{(i)}(x,0)
=
\frac{
K_T^{(i)}(T_0)\theta_L(x,0)
}{
1+
\left[
K_T^{(i)}(T_0)-1
\right]
\theta_L(x,0)
},
\qquad
i=1,2,\ldots,n_t.
\label{eq:ch3_ic_oriani_theta}
\end{equation}

The corresponding initial trapped concentration is

\begin{equation}
C_T^{(i)}(x,0)
=
N_T^{(i)}
\theta_T^{(i)}(x,0).
\label{eq:ch3_ic_oriani_CT}
\end{equation}

The equilibrium constant is defined as

\begin{equation}
K_T^{(i)}(T_0)
=
\exp\left(
\frac{E_b^{(i)}}{RT_0}
\right)
=
\exp\left(
-\frac{\Delta H_i}{RT_0}
\right).
\label{eq:ch3_oriani_KT}
\end{equation}

The initial condition therefore establishes a thermodynamic
partitioning of hydrogen between lattice and trap populations at
$T_0$, while subsequent redistribution during the thermal ramp is
controlled by the finite McNabb--Foster kinetic rates.

%======================================================================
\section{Thermal Desorption}
\label{sec:thermal_desorption}
%======================================================================

\subsection{Programmed Heating Schedule}
\label{subsec:heating_schedule}

During Thermal Desorption Spectroscopy (TDS), the specimen is subjected
to a prescribed temperature ramp. For a constant linear heating rate,
the temperature is expressed as

\begin{equation}
T(t)
=
T_0+\beta t,
\label{eq:ch3_heating_schedule}
\end{equation}

where $T_0$ is the initial desorption temperature and $\beta$ is the
constant heating rate,

\begin{equation}
\beta
=
\frac{dT}{dt}.
\label{eq:ch3_beta}
\end{equation}

The units of $\beta$ may be expressed as
$\mathrm{K\,s^{-1}}$ or $\mathrm{K\,min^{-1}}$, provided that the
corresponding time units are used consistently.

For the present model,

\begin{equation}
T_0=T_{\min},
\label{eq:ch3_Tmin}
\end{equation}

with the initial temperature corresponding to the beginning of the
thermal desorption experiment.

%----------------------------------------------------------------------
\subsection{Thermal Redistribution and Detrapping Kinetics}
\label{subsec:thermal_redistribution}
%----------------------------------------------------------------------

As the temperature increases, the temperature-dependent diffusion,
trapping, and detrapping coefficients change according to their
respective Arrhenius relations.

In particular, the detrapping coefficient increases strongly with
temperature because the detrapping activation energy contains both the
lattice hopping contribution and the trap binding energy,

\begin{equation}
E_d^{(i)}
=
E_L+E_b^{(i)}.
\end{equation}

Consequently, hydrogen initially stored in traps can be released into
the mobile lattice population as the temperature rises
\cite{Turnbull1997,Legrand2015,Marrani2025}.

The released hydrogen increases the local lattice concentration,
producing concentration gradients that drive Fickian diffusion toward
the external surface. Hydrogen reaching the surface is subsequently
released from the specimen and contributes to the measured thermal
desorption signal.

The resulting sequence can therefore be represented as

\begin{equation}
\text{trapped H}
\;\longrightarrow\;
\text{lattice H}
\;\longrightarrow\;
\text{lattice diffusion}
\;\longrightarrow\;
\text{surface}
\;\longrightarrow\;
\text{desorbed H}.
\label{eq:ch3_desorption_sequence}
\end{equation}

The measured TDS response is therefore determined by the coupled
effects of lattice diffusivity, trap density, trap binding energy,
finite trapping/detrapping kinetics, specimen geometry, surface
transport, and heating rate.

%======================================================================
\section{Hydrogen Desorption Flux}
\label{sec:desorption_flux}
%======================================================================

\subsection{Exit Surface Molar Flux}
\label{subsec:exit_flux}

At the external surface $x=L/2$, the instantaneous hydrogen desorption
flux can be obtained from Fick's law:

\begin{equation}
J_{\mathrm{exit}}(t)
=
-\left.
D_L(T(t))
\frac{\partial C_L(x,t)}{\partial x}
\right|_{x=L/2}.
\label{eq:ch3_surface_flux_fick}
\end{equation}

The corresponding units are

\begin{equation}
[J_{\mathrm{exit}}]
=
\mathrm{mol\,m^{-2}\,s^{-1}}.
\end{equation}

When a finite surface activation or desorption resistance is included,
the surface flux may alternatively be represented using a surface
kinetic or penalty relation of the form

\begin{equation}
J_{\mathrm{exit}}(t)
=
\frac{k_{\mathrm{pen}}}{N_L}
\exp\left(
-\frac{E_{\mathrm{bc}}}{RT(t)}
\right)
C_L\left(\frac{L}{2},t\right).
\label{eq:ch3_surface_flux_penalty}
\end{equation}

Here, $k_{\mathrm{pen}}$ and $E_{\mathrm{bc}}$ represent parameters
associated with the adopted surface desorption formulation. The
physical calibration of these parameters must be distinguished from
the bulk diffusion and trapping parameters.

%----------------------------------------------------------------------
\subsection{Conversion of Desorption Flux to Experimental Units}
\label{subsec:flux_conversion}

For comparison with experimentally reported TDS signals, the computed
molar surface flux can be converted into a normalized hydrogen mass
release rate. For a specimen of total thickness $L$, the corresponding
conversion can be written as

\begin{equation}
\mathrm{Flux}_{\mathrm{ppm/s}}(t)
=
\frac{
2M_HJ_{\mathrm{exit}}(t)
}{
L\rho_{\mathrm{steel}}
}
\times10^6,
\label{eq:ch3_flux_ppm_conversion}
\end{equation}

where $M_H$ is the molar mass of atomic hydrogen and
$\rho_{\mathrm{steel}}$ is the mass density of the steel.

For hydrogen,

\begin{equation}
M_H
=
1.008\times10^{-3}
\ \mathrm{kg\,mol^{-1}},
\label{eq:ch3_MH}
\end{equation}

and a representative density of ferritic steel is

\begin{equation}
\rho_{\mathrm{steel}}
\approx
7850\ \mathrm{kg\,m^{-3}}.
\label{eq:ch3_rho}
\end{equation}

The factor of two accounts for desorption from both external surfaces
of a specimen having total thickness $L$, while the numerical model
may solve only the symmetric half-domain $0\leq x\leq L/2$.

%----------------------------------------------------------------------
\subsection{Cumulative Desorbed Hydrogen}
\label{subsec:cumulative_desorption}

The cumulative amount of hydrogen released through the desorption
surface is obtained by integrating the instantaneous flux over time:

\begin{equation}
Q_{\mathrm{des}}
=
\int_0^{t_{\mathrm{final}}}
J_{\mathrm{exit}}(t)\,dt.
\label{eq:ch3_cumulative_desorption_time}
\end{equation}

Using the linear heating relationship

\begin{equation}
dT=\beta\,dt,
\end{equation}

the same quantity can be expressed in temperature coordinates as

\begin{equation}
Q_{\mathrm{des}}
=
\frac{1}{\beta}
\int_{T_0}^{T_{\mathrm{final}}}
J_{\mathrm{exit}}(T)\,dT.
\label{eq:ch3_cumulative_desorption_temperature}
\end{equation}

This relation demonstrates the direct influence of the heating rate on
the conversion between time-dependent hydrogen release and the
temperature-dependent TDS signal.

%======================================================================
\section{Hydrogen Charging Boundary Condition}
\label{sec:charging_boundary}
%======================================================================

\subsection{Mathematical Representations of Hydrogen Entry}
\label{subsec:charging_BC_forms}

Hydrogen charging introduces a surface-entry problem in which hydrogen
must cross the metal--electrolyte interface before becoming part of the
subsurface lattice population. Depending on the physical assumptions
and available experimental information, the charging boundary can be
represented using prescribed concentration, prescribed flux, or a
reaction-controlled surface relation
\cite{Turnbull2014,MartinezPaneda2020,Koren2024a}.

\paragraph{Dirichlet boundary condition.}

The simplest representation prescribes the lattice concentration at
the charging surface:

\begin{equation}
C_L(x_{\mathrm{surface}},t)
=
C_0.
\label{eq:ch3_dirichlet_BC}
\end{equation}

Here, $C_0$ represents the imposed subsurface lattice concentration.
This representation implicitly assumes that the surface-entry process
is sufficiently rapid that the prescribed concentration can be
maintained at the boundary.

\paragraph{Neumann boundary condition.}

Alternatively, the hydrogen entry rate can be represented through a
prescribed hydrogen molar flux:

\begin{equation}
-D_L
\left.
\frac{\partial C_L}{\partial x}
\right|_{\mathrm{surface}}
=
J_H.
\label{eq:ch3_neumann_general}
\end{equation}

When an electrochemical charging current density is related to hydrogen
entry through an effective Faradaic efficiency,
the flux may be written as

\begin{equation}
J_H
=
\frac{
\eta_{\mathrm{Faraday}}\,i_c
}{
F
},
\label{eq:ch3_neumann_electrochemical}
\end{equation}

where $i_c$ is the applied cathodic current density,
$\eta_{\mathrm{Faraday}}$ is the fraction of the electrochemical
hydrogen-generation current associated with hydrogen entering the
metal, and

\begin{equation}
F=96485\ \mathrm{C\,mol^{-1}}
\end{equation}

is Faraday's constant.

Importantly, the applied current density should not automatically be
identified with the hydrogen absorption flux. The conversion depends
on the interfacial electrochemical reactions and therefore requires an
appropriate efficiency or surface-kinetic description
\cite{Turnbull2014,Koren2024a}.

\paragraph{Reaction-controlled or Robin boundary condition.}

A more physically detailed representation accounts for finite surface
absorption and desorption kinetics:

\begin{equation}
-D_L
\left.
\frac{\partial C_L}{\partial x}
\right|_{\mathrm{surface}}
=
k_{\mathrm{abs}}\theta_{\mathrm{ads}}
\left(
1-\frac{C_{L,\mathrm{surf}}}{N_L}
\right)
-
k_{\mathrm{des}}C_{L,\mathrm{surf}}
\left(
1-\theta_{\mathrm{ads}}
\right),
\label{eq:ch3_robin_BC}
\end{equation}

where

\begin{equation}
C_{L,\mathrm{surf}}
=
C_L(x_{\mathrm{surface}},t)
\end{equation}

and $\theta_{\mathrm{ads}}$ represents the fractional surface
coverage of adsorbed hydrogen.

This formulation explicitly separates the surface-entry kinetics from
the subsequent bulk diffusion and trapping processes.

\paragraph{Surface activation penalty formulation.}

The surface desorption or absorption process may also be represented
using an Arrhenius-type surface activation relation:

\begin{equation}
-D_L
\left.
\frac{\partial C_L}{\partial x}
\right|_{x=L/2}
=
\frac{k_{\mathrm{pen}}}{N_L}
\exp\left(
-\frac{E_{\mathrm{bc}}}{RT(t)}
\right)
C_L(L/2,t).
\label{eq:ch3_penalty_BC}
\end{equation}

This formulation provides a finite surface transport resistance while
retaining the bulk McNabb--Foster transport and trapping equations.

%----------------------------------------------------------------------
\subsection{Applied Electrochemical Current and Actual Hydrogen Entry}
\label{subsec:electrochemical_entry}
%----------------------------------------------------------------------

During cathodic electrochemical charging, the externally applied
current density drives electrochemical reactions at the metal--
electrolyte interface. A simplified representation of the initial
hydrogen-generation step is the Volmer reaction,

\begin{equation}
\mathrm{H^+ + e^- \rightarrow H_{ads}}.
\label{eq:ch3_volmer}
\end{equation}

The adsorbed hydrogen can subsequently either enter the metal lattice
or recombine to form molecular hydrogen. Two commonly considered
recombination pathways are the chemical Tafel reaction,

\begin{equation}
2\mathrm{H}_{\mathrm{ads}}
\rightarrow
\mathrm{H}_2,
\label{eq:ch3_tafel}
\end{equation}

and the electrochemical Heyrovsky reaction,

\begin{equation}
\mathrm{H}_{\mathrm{ads}}
+
\mathrm{H^+}
+
e^-
\rightarrow
\mathrm{H}_2.
\label{eq:ch3_heyrovsky}
\end{equation}

Consequently, the externally applied current density $i_c$ does not,
in general, provide a direct one-to-one measurement of the hydrogen
flux entering the metal lattice. The actual absorption flux depends on
the competition between hydrogen adsorption, recombination, and
absorption into the subsurface lattice
\cite{Turnbull2014,Koren2024a}.

The corresponding effective hydrogen entry flux can therefore be
represented as

\begin{equation}
J_H
=
\frac{\eta_{\mathrm{Faraday}}i_c}{F},
\label{eq:ch3_effective_entry_flux}
\end{equation}

where $\eta_{\mathrm{Faraday}}$ should be regarded as an effective
absorption efficiency rather than automatically assuming
$\eta_{\mathrm{Faraday}}=1$.

Hydrogen recombination poisons, such as thiourea or sulfide-containing
species, can modify the surface reaction kinetics and thereby alter
the fraction of generated hydrogen that enters the metal
\cite{Turnbull2014,Ohaeri2018,Koren2024a,Zha2025}.

%----------------------------------------------------------------------
\subsection{Electrochemical Charging and Gaseous Pressure Charging}
\label{subsec:electrochemical_vs_gas}
%----------------------------------------------------------------------

Hydrogen charging can be introduced either through gaseous hydrogen
exposure or through electrochemical cathodic charging. The two
methods produce hydrogen entry through different interfacial
mechanisms and therefore should not be treated as intrinsically
identical boundary conditions.

For gaseous hydrogen exposure, the equilibrium lattice concentration
can be related to hydrogen fugacity through Sieverts-type behaviour:

\begin{equation}
C_0
=
S
\sqrt{f_{\mathrm{H_2}}},
\label{eq:ch3_sieverts}
\end{equation}

where $S$ is the appropriate solubility coefficient and
$f_{\mathrm{H_2}}$ is the hydrogen fugacity.

For electrochemical charging, the subsurface hydrogen concentration
depends on the electrochemical and interfacial conditions, including
the cathodic potential, electrolyte composition, pH, surface
coverage, and the presence of recombination poisons
\cite{Turnbull2014,Koren2024a}.

Therefore, an applied electrochemical potential or current density
cannot be directly converted into an equivalent gas pressure without
an experimentally supported calibration.

One possible bridge between the two charging methods is the
equivalent hydrogen fugacity concept. If the subsurface lattice
concentration generated by electrochemical charging can be measured,
an equivalent gas fugacity may be defined through the same
Sieverts-type relationship:

\begin{equation}
C_{0,\mathrm{EC}}
=
S
\sqrt{
f_{\mathrm{H_2}}^{\mathrm{eq}}
},
\label{eq:ch3_equivalent_fugacity}
\end{equation}

which gives

\begin{equation}
f_{\mathrm{H_2}}^{\mathrm{eq}}
=
\left(
\frac{C_{0,\mathrm{EC}}}{S}
\right)^2.
\label{eq:ch3_equivalent_fugacity_expanded}
\end{equation}

This approach establishes the equivalence through the resulting
subsurface hydrogen concentration rather than by assuming that
electrochemical current, potential, or charge can be directly
identified with gas pressure.

For X65 pipeline steel, experimentally calibrated relationships
between electrochemical charging conditions and equivalent hydrogen
fugacity have been reported in the literature
\cite{Koren2023a,Koren2024a}. Such relationships are material-,
electrolyte-, surface-, and charging-condition-specific and should
therefore be treated as experimental calibration data rather than as a
universal conversion.

%======================================================================
\section{Summary of the Physical Model}
\label{sec:physical_model_summary}
%======================================================================

The physical and mathematical formulation developed in this chapter
can be summarized as a coupled diffusion--reaction system. Mobile
lattice hydrogen undergoes Fickian diffusion,

\begin{equation}
J_L=-D_L(T)\frac{\partial C_L}{\partial x},
\end{equation}

while each trap population exchanges hydrogen with the lattice through
the McNabb--Foster kinetic relation,

\begin{equation}
\frac{\partial C_T^{(i)}}{\partial t}
=
k_i(T)
\frac{N_T^{(i)}}{N_L}
C_L
\left(
1-\frac{C_T^{(i)}}{N_T^{(i)}}
\right)
-
p_i(T)C_T^{(i)}
\left(
1-\frac{C_L}{N_L}
\right).
\end{equation}

The lattice and trapped populations are coupled through the total
hydrogen conservation equation,

\begin{equation}
\frac{\partial C_L}{\partial t}
+
\sum_{i=1}^{n_t}
\frac{\partial C_T^{(i)}}{\partial t}
=
\frac{\partial}{\partial x}
\left(
D_L(T)
\frac{\partial C_L}{\partial x}
\right).
\end{equation}

The initial hydrogen partition is established using the Oriani
equilibrium relation at $T_0$, after which the McNabb--Foster kinetic
equations govern the transient evolution during the programmed
temperature ramp. The thermal history,

\begin{equation}
T(t)=T_0+\beta t,
\end{equation}

couples temperature to lattice diffusion and trapping/detrapping
kinetics. Hydrogen released from traps enters the mobile lattice,
diffuses toward the external surface, and produces the measurable
desorption flux.

The resulting framework provides the physical basis for the numerical
finite-element formulation developed in the subsequent chapter. The
spatial discretization, weak-form derivation, temporal discretization,
nonlinear residual and Jacobian construction, numerical integration,
and computational implementation are therefore treated separately in